% - PER STILARE QUESTO DOCUMENTO VANNO INSERITI I DATI.
% - VISTO CHE IL FILE È LUNGO ED INTRICATO PER SEMPLIFICARE HO MESSO DEI COMMENTI
% - CHE INDICANO I PUNTI DOVE INSERIRE LE INFORMAZIONI CHE MANCANO.
% - INSERISCILI IN CORRISPONDENZA DEI COMMENTI:

% - "INSERT TITOLO"
% - "INSERT USO"
% - "INSERT DATA"
% - "INSERT DESCRIZIONE"
% - "INSERT INFORMAZIONI INCONTRO"
% - "INSERT DOMANDE E RISPOSTE" (intanto lo strutturiamo così, con domande e rispose, se servono altri tipi di struttura cambieremo)

% - "INSERT TITOLO" (verosimilmente sarà una cosa tipo "Verbale Esterno/Interno del xx/yy") 
\newcommand{\Titolo}{Verbale riunione 19 Aprile 2018}

\newcommand{\Gruppo}{SWEefty}


\newcommand{\Redazione}{Elia Montecchio}

\newcommand{\ACapoRedazione}{Elia Montecchio}


\newcommand{\Distribuzione}{Prof.Tullio Vardanega \newline Prof.Riccardo Cardin \newline SWEefty \newline IKS S.r.l.}

\newcommand{\InformazioniDocumento}{}
% - "INSERT USO (INTERNO/ESTERNO)"
\newcommand{\Uso}{Esterno}
% - "INSERT DATA"
\newcommand{\Data}{19 Aprile 2018}
% - "INSERT DESCRIZIONE"
\newcommand{\DescrizioneDoc}{Questo documento riporta in modo formale il riassunto della rinuione tenutasi nella data sopra riportata.  }



\documentclass[a4paper, oneside, openany]{article}

% \usepackage{graphbox}
% permette di modificare i margini
\usepackage[top=3.1cm, bottom=3.1cm, left=2.2cm, right=2.2cm]{geometry}

\usepackage{lastpage} %info sul # dell'ultima pagina del documento
\usepackage{fancyhdr} %per modificare dimensioni,margini, intestazioni e righe a piè di pagina
\fancypagestyle{plain}{
  % cancella tutti i campi di intestazione e piè di pagina
  \fancyhf{}
  
  \lhead{\includegraphics[width=3cm]{../../../CommonImages/logo.jpg}}
  
  \rhead{sweeftyteam@gmail.com}
  
  \lfoot{ %piè di pagina
   {\Titolo} \ - \textit{{\Gruppo}}
  }
  \rfoot{Pagina \thepage{} di \pageref{LastPage}} %es: pag: 4 di 10

  %linea orizzontale alle posizioni top e bottom della pagina
  \renewcommand{\headrulewidth}{0.2	pt}  
  \renewcommand{\footrulewidth}{0.2pt}
}
\pagestyle{plain}

%Comando Spazio
\newcommand{\Spazio}{\mbox{} \\ \mbox{} \\ }  


%\usepackage{calc} %introduce la notazione infissa per le op. aritmetiche interne a LaTeX

\usepackage[utf8]{inputenc}
\usepackage[T1]{fontenc}
\usepackage[italian]{babel} %il documento è in italiano
%\usepackage{textcomp} %The pack­age sup­ports the Text Com­pan­ion fonts, which pro­vide many text sym­bols
%(such as baht, bul­let, copy­right, mu­si­cal­note, onequar­ter, sec­tion, and yen), in the TS1 en­cod­ing.

\usepackage{graphicx}       %permette di inserire delle immagini
\usepackage{caption}        %numerazione figure e loro descrizione testuale
\usepackage{subcaption}     %sottofigure numerabili
\usepackage{float}  %permette di inserire un # qualsiasi di figure fluttuanti
\usepackage{xcolor}
\usepackage{rotating} %permette di ruotare le immagini
%\usepackage{changepage} %utile se c'è bisogno di aggiustare margini per centrare figure

%package utili per la math mode ( $ ... $ o \[ ... \] )
\usepackage{amsmath}
\usepackage{amssymb}
\usepackage{amsfonts}
%\usepackage{euler}    %font 'ams euler', lo stesso di 'Concrete Mathematics' di Knuth
\usepackage{amsthm}
\usepackage{mathtools}

% package utili per tabelle(\thead in particolare)
\usepackage{array, booktabs, caption}
\usepackage{makecell}
\renewcommand\theadfont{\bfseries}
\usepackage{boldline}

\usepackage{listings} %permette di inserire degli spezzoni di codice

% \usepackage{tikz} %disegno di immagini vettoriali a schermo. Utile per grafi
% \usetikzlibrary{arrows.meta}
% \usetikzlibrary{graphs}
% \usetikzlibrary{arrows}
%\usepackage{tikz-uml} %serve per disgnare l'UML, fantastica guida:
%https://perso.ensta-paristech.fr/~kielbasi/tikzuml/var/files/doc/tikzumlmanual.pdf
%download package: http://perso.ensta-paristech.fr/~kielbasi/tikzuml/

%package per le tabelle
\usepackage{booktabs} %permette di poter usare delle liste nelle tabelle
\usepackage{tabularx} 
\usepackage{longtable} %una tabella può continuare su più pagine
\usepackage{multirow} %utile per visualizzare una cella su più righe
%\usepackage{multicolumn} %cella su più colonne
%\usepackage[table]{xcolor} %rende disponibile l'utilizzo di un colore per lo sfondo
                        %delle celle di una tabella

%crea una cella per le tabelle in grado di andare a capo con \newline
%https://tex.stackexchange.com/questions/12703/how-to-create-fixed-width-table-columns-with-text-raggedright-centered-raggedlef
\usepackage{array}
\newcolumntype{L}[1]{>{\raggedright\let\newline\\\arraybackslash\hspace{0pt}}m{#1}}
\newcolumntype{C}[1]{>{\centering\let\newline\\\arraybackslash\hspace{0pt}}m{#1}}
\newcolumntype{R}[1]{>{\raggedleft\let\newline\\\arraybackslash\hspace{0pt}}m{#1}}


%indice con i puntini
\usepackage{tocloft}
\renewcommand\cftsecleader{\cftdotfill{\cftdotsep}}

%http://ctan.mirror.garr.it/mirrors/CTAN/macros/latex/contrib/appendix/appendix.pdf
\usepackage{appendix} %aggiunge dei comandi per l'appendice
\usepackage{parskip} %aiuta LaTeX a trovare il miglior stile per i page break
\setcounter{secnumdepth}{5} % numera i sottoparagrafi
\setcounter{tocdepth}{5} %aggiunge all'indice i sottoparagrafi
%\usepackage{titlesec} %\begin{paragraph} si può usare come subsubsubsection!


\usepackage{breakurl}%\url{...} può continare alla linea successiva. (si può andare a capo)

\definecolor{Maroon}{cmyk}{0, 0.87, 0.68, 0.32}
\usepackage[colorlinks=true]{hyperref}
\hypersetup{
    colorlinks=true,
    citecolor=black,
    filecolor=black,
    linkcolor=black, % colore dei link interni
    urlcolor=Maroon  % colore dei link interniesterni
}

%impostazioni per il codice che deve finire dentro a
%\begin{lstlisting}

\definecolor{listinggray}{gray}{0.9}
\definecolor{lbcolor}{rgb}{0.9,0.9,0.9}
\lstset{
backgroundcolor=\color{lbcolor},
    tabsize=4,    
%   rulecolor=,
    language=[GNU]C++,
    basicstyle=\scriptsize,
    upquote=true,
    aboveskip={1.5\baselineskip},
    columns=fixed,
    showstringspaces=false,
    extendedchars=true,
    inputencoding=utf8,
    breaklines=true,
    prebreak = \raisebox{0ex}[0ex][0ex]{\ensuremath{\hookleftarrow}},
    frame=single,
    numbers=left,
    showtabs=false,
    showspaces=false,
    showstringspaces=false,
    identifierstyle=\ttfamily,
    keywordstyle=\color[rgb]{0,0,1},
    commentstyle=\color[rgb]{0.026,0.112,0.095},
    stringstyle=\color[rgb]{0.627,0.126,0.941},
    numberstyle=\color[rgb]{0.205, 0.142, 0.73},
%        \lstdefinestyle{C++}{language=C++,style=numbers}’.
}
\lstset{
  backgroundcolor=\color{lbcolor},
  tabsize=4,
  language=C++,
  captionpos=b,
  tabsize=3,
  frame=lines,
  numbers=left,
  numberstyle=\tiny,
  numbersep=5pt,
  breaklines=true,
  showstringspaces=false,
  basicstyle=\footnotesize,
  identifierstyle=\color{magenta},
  keywordstyle=\color[rgb]{0,0,1},
  commentstyle=\color{orange},
  stringstyle=\color{red}
}


 \newgeometry{top=4cm}

\begin{document}
\begin{titlepage}
	\begin{center}
		
		\begin{center}
			%% qui metteteci l'immagine di copertina. Io ho messo quella dell'uni,
			%voi mettete quella del vostro grupo
			\centerline{\includegraphics[scale=0.24]{../../../CommonImages/logo.jpg}}
		\end{center}
		
		\vspace{1cm}
		
		\begin{Huge}
			\textbf{\Titolo{}} \\
		\end{Huge}
		
		\vspace{9pt}  
		
		\begin{large}
			\Gruppo{}\ - \Data{}
		\end{large}	  
		
		\vspace{15pt}
		
		\bgroup
		\def\arraystretch{1.3}
		\centering
		\begin{tabular}{c|L{5cm}}
			\multicolumn{2}{c}{\textbf{\InformazioniDocumento{}} } \\ \hline
			Redazione & \ACapoRedazione{} \\

			Uso & \Uso \\
			Distribuzione & \Distribuzione{}
		\end{tabular}
		\egroup
		
		\vspace{15pt}
		
		\begin{center}
			\textbf{Descrizione\\}
			\DescrizioneDoc{}
		\end{center}
		
	\end{center}
\end{titlepage}

\restoregeometry
	
	\section{Informazioni generali}
		\subsection{Informazioni incontro}
			% - INSERT INFORMAZIONI INCONTRO
			\begin{itemize}
				\item { \textbf{Luogo:} Torre Archimede  }
				\item { \textbf{Data:} 19 Aprile 2018 }
				\item { \textbf{Ora inizio:} 17.30 }
				\item { \textbf{Ora fine:} 17.50 }
				\item { \textbf{Partecipanti del gruppo:} Gruppo al completo }
				\item { \textbf{Partecipanti esterni:} IKS Srl }
			\end{itemize}
		
	
	\subsection{Argomenti affrontati}
    Nel corso della riunione effettuata mediante Skype, è stato mostrato all'azienda lo stato dei plugin. Inoltre, sono state discusse varie problematiche emerse durante lo sviluppo del software.
    

 
	 Sono stati illustrati alla proponente i vari incrementi effettuati riguardanti la stack trace e il plugin della mappa topologica. 
	I due plugin attualmente si presentano in un'unica pagina ma in seguito verranno divisi.
	Il proponente come prima cosa chiede se la mappa topologica è interagibile e se nel caso vi siano più nodi del sistema, rispetto a quelli attualmente visibili, vengono visualizzati nella mappa.
	Per questo motivo viene mostrato che è possibile effettuare delle operazioni nella mappa interagendo direttamente con essa.
	Nello specifico è possibile ingrandire e rimpicciolire la mappa, spostare ciascun nodo visualizzato e spostare interamente la mappa.
	Inoltre, la mappa è dinamica dunque nel caso venissero aggiunti nodi al sistema, importando appositi documenti in ElasticSearch, il plugin è in grado di mapparli. 
	In seguito viene spiegato al proponente che il tempo media di risposta che viene visualizzato nella linea che unisce due nodi del sistema è anch'esso dinamico e viene calcolato in un modo specifico.
	
	Successivamente il team si confronta con la proponente riguardo la rappresentazione di un cluster del sistema all'interno della mappa topologica in quanto risulta essere un requisito obbligatorio ancora da implementare. Il team, infatti, non riesce a capire quale tipo di documenti deve prelevare dal database per rappresentare questo tipo di nodo.
	Assieme alla proponente si è cercato di sviscerare il problema andando effettivamente a controllare i vari documenti presenti all'interno di ElasticSearch effettuando delle query.   
	In seguito ad una verifica accurata della proponente è emerso che non è presente alcuna informazione che riporti ad un nodo che rappresenti un cluster all'interno del sistema, in quanto il sistema che si sta monitorando attualmente registra solo chiamate da client a server e non vi sono cluster.
	Per questo motivo si è deciso \textbf{01/VE\_2018-04-19} di spostare un requisito da obbligatorio a opzionale.
	
	La proponente in seguito chiede al team di poter visionare lo stato e il funzionamento del secondo plugin: la stack trace.
	Viene illustrata la stack trace che si presenta come una tabella, dove ogni riga rappresenta una specifica richiesta al server dal client. Ogni richiesta ha un nome, un tempo di esecuzione e indica se durante l'esecuzione si è presentato un errore. Cliccando su ogni richiesta, vi è la possibile di aprire un popup nel quale è visibile un dettaglio di informazioni specifiche della trace. 
	Nello specifico vi è il call tree, se presente, con il self time e total time di ogni funzione. Si può accedere anche alla query list dove sono indicate tutte le query effettuate nel database al fine di fornire una risposta al client.
	La proponente chiede se i dati di ogni richiesta sono relativi ad un solo documento memorizzato nel database oppure da un raggruppamento di documenti.
	Per questo motivo viene spiegato che le richieste HTTP risultano essere singole, quelle di PAGE LOAD sono la combinazione tra una richiesta HTTP e quelle con lo stesso trace id che permette di aggiungere la durata di caricamento della pagina.
	Mentre per le richiste JDBC vengono aggiunte le trace con lo stesso child span riportato nelle richieste HTTP e di PAGE LOAD.
	Dunque a seconda del tipo di richiesta, ogni riga ha informazioni provenienti da diversi documenti.
	
	La proponente risulta essere soddisfatta del lavoro svolto finora e
	dopo un confronto si è deciso di procedere nel seguente modo:
	    \begin{itemize}
	    \item { \textbf{02/VE\_2018-04-19} Il team dovrà inviare alla proponente degli screenshot inerenti al lavoro svolto; }
	    \item { \textbf{03/VE\_2018-04-19} Il team dovrà inviare alla proponente una mail presentando il link alla repository github che contiene il codice del prodotto. }	
	    	
	    \end{itemize}
	
	
	
	
	
	
	

		

\end{document}